\documentclass[10pt,a4paper]{article}
\usepackage{amsmath,amssymb,amsthm,graphicx}
\usepackage{amsmath,amssymb,amsthm}
\newtheorem{theorem}{Theorem}
\newtheorem{proposition}[theorem]{Proposition}
\newtheorem{claim}[theorem]{Claim}
\newcommand{\norm}[1]{\lVert#1\rVert}
\title{An explicit three-player NOF exactly-N protocol with improved leading communication constant}
\author{Nati Linial, Adi Shraibman}
\date{doi:10.4230/LIPIcs.CCC.2021.2; CC BY 4.0}
\begin{document}
\maketitle
\noindent\textit{Editorial note.} Editable transcription of the official CC BY 4.0 eight-page publication. Complete original Section2 proof retained; model and main theorem from Section1, direct Section3 consequence retained as context. Original entropy formulas, carry inequalities, cover weights and q/d choices are transcribed without mathematical correction or added assumptions. Standard background historical discussion and conjectures omitted. Bracket citation numbers retain original 23-entry bibliography numbering. Companion arxiv/journal identities denote the same result and are not counted again. All production and validation on Yupeng; compilation belongs to root.\par\medskip
\noindent\textit{Allowed standard prerequisites (editorial).}\begin{itemize}
\item NOF blackboard communication, conditional Shannon entropy, standard integral versus fractional set-cover bound. Official Sections1–2; standard Lovasz set-cover theorem cited as [15] with complete source-local use included
\end{itemize}\medskip
\section*{Model and conventions (original Section 1)}
The multiplayer Number On the Forehead (NOF) model of communication complexity was introduced by Chandra, Furst and Lipton [9]. Given a function $f:[N]^k\to\{0,1\}$, the $k$ players in this scenario should jointly find out $f(x_1,\ldots,x_k)$. We think of $x_i$ as being placed on player $i$'s forehead, so that each player sees the whole input bar one argument. Players communicate by writing bits on a shared blackboard according to an agreed-upon protocol.
The exactly-$N$ function $f:[N]^k\to\{0,1\}$ is given by $f(x_1,\ldots,x_k)=1$ iff $\sum x_i=N$.
All logarithms in this paper are in base 2. Also, unless otherwise specified, all asymptotic statements are taken with $N\to\infty$.
\begin{theorem}\label{main}
There is an explicit protocol for NOF exactly-$N$ of complexity
\begin{equation}\label{eq:main}
2\sqrt{\log e}\sqrt{\log N}+o(\sqrt{\log N})=2.4022\ldots\sqrt{\log N}+o(\sqrt{\log N}).\tag{2}
\end{equation}
\end{theorem}

\section*{Proof of Theorem 1 (original Section 2)}
The three players in our protocol are called $P_x$, $P_y$ and $P_z$. The inputs that they get to see are $(y,z)$, $(x,z)$ and $(x,y)$ respectively.
Here is a similar problem in the realm of vector addition. Given integers $q,d>1$, define $g=g_{q,d}(\alpha,\beta,\gamma)$ to be 1 if $\alpha+\beta=\gamma$ and 0 otherwise. Here $\alpha,\beta\in[q]^d$, $\gamma\in[2q]^d$ and addition is vector addition in $\mathbb R^d$. The following one-round protocol [3] for $g$ is correct because the inequality $\norm{2\alpha-\gamma}^2+\norm{2\beta-\gamma}^2\ge2\norm{\alpha-\beta}^2$ holds always and is an equality iff $\gamma=\alpha+\beta$.
\paragraph{Protocol 1. A protocol for $g_{q,d}$.}
\begin{enumerate}
\item $P_z$ computes $\norm{\alpha-\beta}_2^2$, and writes the result on the board.
\item $P_y$ writes 1 iff $\norm{\alpha-\beta}_2^2=\norm{2\alpha-\gamma}_2^2$.
\item $P_x$ writes 1 iff $\norm{\alpha-\beta}_2^2=\norm{2\beta-\gamma}_2^2$.
\end{enumerate}
The cost of this protocol is $2+\log dq^2$.
The above is an efficient method to decide high-dimensional vector addition, but our objective is to decide the integer addition relation $X+Y+Z=N$. We let $x=X$, $y=Y$ and $z=N-Z$, so the relation we need to consider is $x+y=z$.
Our protocol to decide whether $x+y=z$ builds on the protocol for $g_{q,d}$. It is the issue of carry bits in integer addition that makes this decision problem harder. The integers $q,d>1$ are chosen so that
\begin{equation}2qN>q^d\ge2N.\tag{3}\label{eq:range}\end{equation}
the specific choice is made below so as to minimize the cost of the protocol.
We denote by $w_q$ the vector that corresponds to the base $q$ representation of the integer $w$. As usual, $e_i$ is the $d$-dimensional vector with 1 in the $i$-th coordinate and zeros elsewhere.
Let $C(x,y)\in\{0,1\}^d$ be the carry vector when $x$ and $y$ are added in base $q$. The relation $x+y=z$ among integers is equivalent to the vector relation
\[x_q+y_q=\zeta,\]
where the $i$-th coordinate of $\zeta$ is
\[\zeta_i=z_i+q\cdot C(x,y)_i-C(x,y)_{i-1}\]
(Here $C(x,y)_0=0$).
The protocol from [3] now suggests itself: $P_z$ posts $C(x,y)$, and Protocol 1 is used to decide the relation $x_q+y_q=\zeta$. This yields again the estimate (1) of the original paper, $2\sqrt2\sqrt{\log N}+o(\sqrt{\log N})$.
The alternative approach that we adopt here considers instead the equivalent vector relation
\[x_q+\eta=z_q\]
where
\[\eta=(x+y)_q-x_q.\]
Concretely, the $i$-th coordinate of $\eta$ is:
\[\eta_i=y_i-q\cdot C(x,y)_i+C(x,y)_{i-1}.\]
In order to run Protocol 1, $P_z$ needs to know $\eta$ and $x_q$, which he does. The situation with $P_y$ is even simpler, since he needs to know $x_q$ and $z_q$ which are his inputs. The only difficulty is with $P_x$ who needs to know $z_q$ (which he does) and $\eta$. The latter is not part of his input and $P_z$ fills in the missing information for him.
The obvious solution is for $P_z$ to reveal $C$ to $P_x$ using $d$ bits of information. However, we can save communication by exploiting the fact that $P_x$ and $P_z$ share some information, i.e., they both know $y$ for every $y\ne0$.
By a standard argument in this area which we detail below (Proposition 2), a protocol that works for typical pairs $x,y$ can be easily modified to work in all cases. So, let us pick $x$ and $y$ uniformly at random from among the $d$-digit numbers in base $q$ and think of $C$, the vector of carry bits as a random variable on this probability space. The number of bits that $P_z$ needs to post so that $P_x$ gets to know $C$, and therefore know $\eta$, is $H(C\mid y)$, the entropy of $C$ given $y$. The gain is clear, since $H(C)\ge H(C\mid y)$.
It remains to estimate $H(C\mid y)$. Fix some integers $s\ge t\ge0$, and let $X$ be the random variable that is a uniformly sampled subset of $[s]$ of cardinality $\ge t$. It is easily verified that $H(X)=(1+o_s(1))\cdot s\cdot h(t/s)$, where $h(\cdot)$ is the univariate entropy function. The entropy of $X$ is the same also if we sample subsets of $[s]$ of cardinality $<t$. Let $r$ be an integer in the range $d\gg r\gg1$, e.g., $r\approx\sqrt d$. For $j=1,\ldots,r$, let
\[S_j=\{i\mid \tfrac{qj}{r}>y_i\ge\tfrac{q(j-1)}{r}\},\]
where $q>y_i\ge0$ is the $i$-th digit of $y$. A carry occurs in digit $i\in S_j$ only if $x_i>\frac{q(r-j)}r$, where $x_i$ is the $i$-th digit of $x$. Then
\[H(C\mid y)\le(1+o_r(1))\sum_{j=1}^r\frac{|S_j|}{d}h(\tfrac jr).\]
Since $y$ is chosen at random, $|S_j|\le(1+o_r(1))\frac dr$, and so
\[H(C\mid y)\le(1+o_r(1))\sum_{j=1}^r\frac1r h(\tfrac jr).\]
The limit of this expression as $r\to\infty$ is
\[\lambda=\int_0^1h(u)\,du=\frac{\log e}{2}=0.721\ldots\]
It is left to optimize on $q$ and $d$. The complexity of our protocol is
\[\lambda d+\log dq^2+2,\]
where recall that $2qN>q^d\ge2N$. It is not hard to verify that choosing
\begin{equation}d=\sqrt{\frac2\lambda\log2N}\qquad q=2^{\sqrt{\frac\lambda2\log2N}},\tag{4}\label{eq:choice}\end{equation}
we get a protocol with complexity
\[2\sqrt{2\lambda\log N}+o(\sqrt{\log N}),\]
and this is asymptotically optimal in our setting.
To sum up, here is the protocol which proves Theorem 1:
\paragraph{Protocol 2. A protocol for exactly-$N$, for typical pairs $x,y$.}
For $d,q$ as in Equation (4)
\begin{enumerate}
\item $P_z$ publishes the vector $\eta=(x+y)_q-x_q$ in a way that $P_x$ can read.
\item The players run protocol 1 for $g_{q,d}$ on inputs $x_q,\eta,z_q$. That is:
\begin{enumerate}
\item $P_z$ writes $\norm{\eta-x_q}_2^2$ on the board
\item $P_y$ writes 1 iff $\norm{\eta-x_q}_2^2=\norm{2x_q-z_q}_2^2$.
\item $P_x$ writes 1 iff $\norm{\eta-x_q}_2^2=\norm{2\eta-z_q}_2^2$.
\end{enumerate}
\end{enumerate}
\begin{proposition}
Let $P$ be an NOF protocol for the exactly-$N$ that works correctly for an $\Omega(1)$-fraction of the input pairs $x,y$ (and every $z$) with communication complexity $\Phi(N)$. Then there is an NOF protocol that works for all inputs with communication complexity $\Phi(N)+O(\log\log N)$.
\end{proposition}
\begin{proof}
Let $S\subseteq[N]\times[N]$ be the set of input pairs $x,y$ on which $P$ succeeds. We claim that there is a collection $F$ of $O(\log N)$ vectors $\Delta\in[N]\times[N]$ such that
\[\bigcup_{\Delta\in F}(S+\Delta)\supseteq[N]\times[N].\]
In the modified protocol, $P_z$ sees $x,y$ and announces the index of some $\Delta=(\Delta_1,\Delta_2)\in F$ for which $(x-\Delta_1,y-\Delta_2)\in S$. Then the players run Protocol 2 with inputs $(x-\Delta_1,y-\Delta_2,z-\Delta_1-\Delta_2)$.
The construction of $F$ uses a standard fact about the set-cover problem. For a family of finite sets $\mathcal X\subseteq2^\Omega$ we denote by $c(\mathcal X)$ the least number of members in $\mathcal X$ whose union is $\Omega$. Also $c^*(\mathcal X)$ is the minimum cost of a fractional cover. Namely,
\[c^*(\mathcal X)=\min\sum_{\mathcal X}\omega_X,\quad\text{where }\omega_X\ge0\text{ for every }X\in\mathcal X\text{ and }\sum_{x\in X}\omega_X\ge1\text{ for every }x\in\Omega.\]
Then
\[c(\mathcal X)\le\log(|\Omega|)\cdot c^*(\mathcal X)\]
(e.g., Lov\'asz [15]) and actually the greedy algorithm yields a set cover that meets this bound.
In our case $\Omega=[N]\times[N]$, and
\[\mathcal X=\{(S+\Delta)\cap([N]\times[N])\mid\Delta\in[-N,N]\times[-N,N]\}.\]
It is easily verified that the weights $\omega_x=\frac{10}{N^2}$ constitute a fractional cover, so that $c^*(\mathcal X)\le40$ and hence $c(\mathcal X)\le80\log N$, as claimed.
\end{proof}

\section*{Applications in additive combinatorics (original Section 3, contextual consequence)}
Let $\rho_k(N)$ be the largest density of a subset of $[N]$ without an arithmetic progression of length $k$. A corner is a triple of points $(x,y),(x+\delta,y),(x,y+\delta)$ for some $\delta\ne0$. Denote the largest density of a subset of $[N]\times[N]$ without a corner by $\rho_3^\angle(N)$.
As mentioned in the introduction, the NOF communication complexity of $f$ is closely related to corners theorems. Our Theorem 1 immediately implies,
\begin{theorem}
\[\rho_3^\angle(N)\ge2^{-2\sqrt{\log e}\sqrt{\log N}+o(\sqrt{\log N})}=2^{-2.4022\ldots\sqrt{\log N}+o(\sqrt{\log N})}.\]
There is an explicit corner-free subset of $[N]\times[N]$ of size
\[N^2/2^{2\sqrt{\log e}\sqrt{\log N}+o(\sqrt{\log N})}.\]
\end{theorem}
The derivation of Theorem 3 from Theorem 1 is an easy consequence of the following claim.
\begin{claim}[{[9], implicit}]
\begin{enumerate}
\item There is an optimal one-round protocol for the addition problem.
\item Let $T=T(x,y)$ be the message that the $P_z$ sends on inputs $(x,y)$ in a one-round protocol for the addition problem. Then the set
\[S(T)=\{(x,y):T(x,y)=T\}\]
is corner-free.
\end{enumerate}
\end{claim}
See [9,7,1,14,3] for more details about the above claim and the relation between communication complexity and additive combinatorics. The same comments and corollaries above apply also verbatim to the $(6,3)$ Theorem (e.g., [19]) and to the quantitative version of the triangle removal lemma.

\begin{thebibliography}{23}
\bibitem{ref1} A. Ada, A. Chattopadhyay, O. Fawzi, and P. Nguyen. The NOF multiparty communication complexity of composed functions. computational complexity, 24(3):645–694, 2015.
\bibitem{ref2} M. Ajtai and E. Szemerédi. Sets of lattice points that form no squares. Stud. Sci. Math. Hungar, 9(1975):9–11, 1974.
\bibitem{ref3} N. Alon and A. Shraibman. Number on the forehead protocols yielding dense ruzsa–szemerédi graphs and hypergraphs. Acta Mathematica Hungarica, 161(2):488–506, 2020.
\bibitem{ref4} P. Beame, M. David, T. Pitassi, and P. Woelfel. Separating deterministic from randomized nof multiparty communication complexity. In Proceedings of the 34th International Colloquium On Automata, Languages and Programming, Lecture Notes in Computer Science. Springer-Verlag, 2007.
\bibitem{ref5} P. Beame, T. Pitassi, and N. Segerlind. Lower bounds for Lovász-Schrijver systems and beyond follow from multiparty communication complexity. SIAM Journal on Computing, 37(3):845–869, 2006.
\bibitem{ref6} F. A. Behrend. On sets of integers which contain no three terms in arithmetical progression. Proceedings of the National Academy of Sciences, 32(12):331–332, 1946.
\bibitem{ref7} R. Beigel, W. Gasarch, and J. Glenn. The multiparty communication complexity of Exact-T: Improved bounds and new problems. In International Symposium on Mathematical Foundations of Computer Science, pages 146–156. Springer, 2006.
\bibitem{ref8} T. F. Bloom and O. Sisask. Breaking the logarithmic barrier in Roth’s theorem on arithmetic progressions. arXiv preprint, 2020. arXiv:2007.03528.
\bibitem{ref9} A. Chandra, M. Furst, and R. Lipton. Multi-party protocols. In Proceedings of the 15th ACM Symposium on the Theory of Computing, pages 94–99. ACM, 1983.
\bibitem{ref10} M. Elkin. An improved construction of progression-free sets. In Proceedings of the twenty-first annual ACM-SIAM symposium on Discrete Algorithms, pages 886–905. Society for Industrial and Applied Mathematics, 2010.
\bibitem{ref11} W. T. Gowers. Hypergraph regularity and the multidimensional szemerédi theorem. Annals of Mathematics, pages 897–946, 2007.
\bibitem{ref12} J. Håstad and M. Goldmann. On the power of small-depth threshold circuits. Computational Complexity, 1:113–129, 1991.
\bibitem{ref13} E. Kushilevitz and N. Nisan. Communication Complexity. Cambridge University Press, 1997.
\bibitem{ref14} N. Linial, T. Pitassi, and A. Shraibman. On the communication complexity of high-dimensional permutations. In 10th Innovations in Theoretical Computer Science Conference, ITCS San Diego, California, USA, volume 124, pages 54:1–54:20, 2019.
\bibitem{ref15} L. Lovász. On the ratio of optimal integral and fractional covers. Discrete Mathematics, 13:383–390, 1975.
\bibitem{ref16} B. Nagle, V. Rödl, and M. Schacht. The counting lemma for regular k-uniform hypergraphs. Random Structures \& Algorithms, 28(2):113–179, 2006.
\bibitem{ref17} V. Rödl and J. Skokan. Regularity lemma for k-uniform hypergraphs. Random Structures \& Algorithms, 25(1):1–42, 2004.
\bibitem{ref18} K. F. Roth. On certain sets of integers. Journal of the London Mathematical Society, 1(1):104– 109, 1953.
\bibitem{ref19} I. Ruzsa and E. Szemerédi. Triple systems with no six points carrying three triangles. Combinatorics (Keszthely, 1976), Coll. Math. Soc. J. Bolyai, 18:939–945, 1978.
\bibitem{ref20} J. Solymosi. Note on a generalization of Roth’s theorem. Discrete and Computational Geometry: The Goodman-Pollack Festschrift, pages 825–827, 2003.
\bibitem{ref21} E. Szemerédi. On sets of integers containing no k elements in arithmetic progression. Acta Arith, 27(199-245):2, 1975.
\bibitem{ref22} B. L. van der Waerden. Beweis einer Baudetschen Vermutung. Nieuw Arch. Wiskunde, 15:212–216, 1927.
\bibitem{ref23} A. Yao. On ACC and threshold circuits. In Proceedings of the 31st IEEE Symposium on Foundations of Computer Science, pages 619–627. IEEE, 1990.
\end{thebibliography}
\end{document}
